% ============================================================
%  SECTION 2 — LITERATURE REVIEW
% ============================================================

\subsection{Overview and Review Strategy}
\label{subsec:lr_overview}
% ------------------------------------------------------------

This section maps the intellectual lineage of the present work across four
adjacent knowledge domains: project management~(PM) foundations, project
portfolio management~(PPM) and selection, PPM budgeting and cash-flow
modelling, and reinforcement learning~(RL) applied to sequential
resource-allocation problems. Together these streams define both the
theoretical foundation and the gap that motivates the proposed framework.
The review follows the PRISMA~2020 reporting guidelines \citep{Page2021}
to ensure transparency and reproducibility in source identification,
screening, and inclusion decisions.

% - - - - - - - - - - - - - - - - - - - - - - - - - - - - - -
\subsubsection{Search and Screening Protocol}
\label{subsubsec:prisma}
% - - - - - - - - - - - - - - - - - - - - - - - - - - - - - -

Literature was retrieved from Google Scholar, Semantic Scholar, and Scopus
using a progressively narrowing search strategy. Queries began at the
broadest relevant domain (\emph{project management}) and were successively
refined through \emph{project portfolio management}, \emph{portfolio
budgeting and cash-flow optimisation}, \emph{earned value management under
uncertainty}, and \emph{reinforcement learning for project scheduling and
resource allocation}. Citation-network traversal --- backward and forward
snowballing from high-citation anchor papers --- was used to recover seminal
works not surfaced by keyword search. For instance, a high-citation anchor
paper on portfolio optimisation exposes its own foundational references
through backward traversal, while forward traversal through its citing works
surfaces more recent developments; both layers would have been invisible to
keyword search alone.

Records were screened in three sequential automated stages consistent with
the PRISMA~2020 flow diagram (Appendix~\ref{app:prisma}). The screening
pipeline, together with its archived corpus and output artefacts, is
implemented as a reproducible processing chain whose repository location
and concept-level mapping are documented in
Appendix~\ref{app:repo_map}.

\medskip
\textbf{Stage~1 --- Identification and Machine Readability.}\;
Candidate PDFs were subjected to automated bibliographic and readability
triage. Each document was assessed for text-extraction yield (minimum
readable character density across pages), lexical quality
(alphabetic-token ratio and mean word length), and sentence-transformer
embeddability. Documents failing any criterion were quarantined in a
dedicated rejection archive, with a content hash logged to prevent
duplicate submissions from re-entering the pipeline.

\medskip
\textbf{Stage~2 --- Relevance Filtering.}\;
Records passing Stage~1 were ranked by a composite relevance score. A
semantic similarity component embedded each paper's abstract,
introduction, and conclusion using a sentence-transformer model and
computed the maximum cosine similarity against twelve thesis-derived
query vectors spanning the target knowledge domains. A keyword-scoring
component assigned weighted counts of domain-specific high-value and
medium-value terms, with a penalty for domain-irrelevant vocabulary. The
two components were combined with equal weight ($0.5$ each) to yield a
combined relevance score $r \in [0,1]$. Papers were then triaged into
three cohorts: \emph{highly relevant} ($r \geq 0.45$), \emph{valuable}
($r \geq 0.30$), and \emph{irrelevant} ($r < 0.30$), with the first two
cohorts retained for structured extraction.

\medskip
\textbf{Stage~3 --- Structured Extraction.}\;
Each retained paper in the \emph{highly relevant} and \emph{valuable}
cohorts was subjected to structured knowledge extraction. The abstract,
introduction, and conclusion of each PDF were parsed and submitted to a
large language model with a fixed schema prompt, producing a per-paper
record covering: problem framing, portfolio level, decision type,
stochastic cash-flow treatment, contract mechanics,
termination-settlement provisions, modelling paradigm, calibration
parameters (CPI/SPI statistics, cost-overrun distributions, S-curve shape
parameters), and thesis-relevance flags. The resulting structured corpus
subsequently drove an automated synthesis dashboard and a table
generator, producing the gap-analysis matrix, domain-distribution charts,
and the literature map reproduced in Table~\ref{tab:litmap}.

\medskip

The complete pipeline --- processing scripts, the structured record
corpus, dashboard outputs, and synthesis tables --- is archived within
the project repository and constitutes a reproducible, machine-auditable
literature review corpus; its directory-level organisation is detailed in
Appendix~\ref{app:repo_map}. The PRISMA flow diagram quantifying records
at each stage is provided in Appendix~\ref{app:prisma}.

\begin{figure}[H]
    \centering
    \begin{tikzpicture}[
        box/.style      = {rectangle, rounded corners=6pt, thick,
                           text width=7.0cm, align=center,
                           minimum height=0.85cm, inner sep=7pt,
                           font=\small},
        decision/.style = {diamond, thick, aspect=2.5,
                           minimum width=4.5cm, minimum height=0.7cm,
                           align=center, font=\scriptsize, inner sep=2pt},
        rej/.style      = {rectangle, rounded corners=6pt, thick,
                           text width=3.2cm, align=center,
                           minimum height=0.75cm, inner sep=5pt,
                           font=\scriptsize},
    ]
    
    %% --- Row 1 : left -> right ---
    \node[box, flowColorA] (s1)
        {\textbf{Stage 1 --- Identification \& Readability}\\
         Extraction yield \(\cdot\) Lexical quality \(\cdot\) Embeddability};
    
    \node[decision, flowColorGray, right=1.2cm of s1] (d1)
        {All three tests pass?};
    
    %% --- Row 2 : right -> left ---
    \node[box, flowColorB, below=2.4cm of d1] (s2)
        {\textbf{Stage 2 --- Relevance Filtering}\\
         Semantic similarity \(+\) Keyword score \(\rightarrow r \in [0,1]\)};
    
    \node[decision, flowColorGray, left=1.2cm of s2] (d2)
        {\(r \geq 0.30\)?};
    
    %% --- Row 3 : left -> right ---
    \node[box, flowColorC, below=1.7cm of d2] (s3)
        {\textbf{Stage 3 --- Structured Extraction}\\
         JSON extraction: metadata \(\cdot\) calibration \(\cdot\) thesis flags};
    
    \node[box, flowColorD, right=1.2cm of s3] (out)
        {\textbf{Synthesis}\\
         Gap matrix \(\cdot\) Domain charts \(\cdot\) Literature map};
    
    %% --- Rejection nodes (branch upward into prior row gap) ---
    \node[rej, flowColorReject, above=0.55cm of d1] (r1)
        {Rejected\\(failed readability)};
    \node[rej, flowColorReject, above=0.55cm of d2] (r2)
        {Rejected\\(irrelevant)};
    
    %% --- Main snake flow ---
    \draw[arrowMain] (s1) -- (d1);
    \draw[arrowMain] (d1) -- node[right, font=\scriptsize]{Yes} (s2);
    \draw[arrowMain] (s2) -- (d2);
    \draw[arrowMain] (d2) -- node[left, font=\scriptsize]{Yes} (s3);
    \draw[arrowMain] (s3) -- (out);
    
    %% --- Rejection branches ---
    \draw[arrowReject] (d1) -- node[right, font=\scriptsize]{No} (r1);
    \draw[arrowReject] (d2) -- node[left, font=\scriptsize]{No} (r2);
    
    \end{tikzpicture}
    \caption{Three-stage PRISMA-aligned screening pipeline, arranged as a
    serpentine flow. Boxes denote automated processing stages; diamonds
    denote pass/fail gates; dashed arrows indicate rejection routing.
    Implementation details and repository locations are provided in
    Appendix~\ref{app:repo_map}.}
    \label{fig:screening_pipeline}
\end{figure}


% ------------------------------------------------------------
\subsection{Project Management: Foundational Concepts}
\label{subsec:lr_pm}
% ------------------------------------------------------------

The academic study of project management draws on schedule control (CPM/PERT),
earned value management~(EVM), and risk modelling as its core quantitative
pillars \citep{Fleming2010, PMI2021}. EVM introduced the Cost Performance Index
(CPI) and Schedule Performance Index (SPI) as real-time diagnostic signals
relating budgeted cost of work performed to actual expenditure and planned
schedule --- providing the performance vocabulary that the present model
inherits directly. The present model adopts the Earned Schedule correction
to SPI \citep{Lipke2003}, replacing the classical ratio $\mathrm{BCWP}/\mathrm{BCWS}$
with $\mathrm{SPI}(t) = \mathrm{ES}/\mathrm{AT}$, which continues to degrade
past the planned finish date rather than freezing when $\mathrm{BCWS}$ reaches
$\mathrm{BAC}$ --- a known limitation of classical SPI in post-planned-finish
recovery periods \citep{Lipke2003, Henderson2004}. The To-Complete Performance Index
(TCPI) is also included as a forward-looking complement to CPI, measuring the cost
efficiency required on remaining work to finish within budget. Values of $\mathrm{CPI} < 1$ signal cost overrun;
$\mathrm{SPI} < 1$ signals schedule lag. Stochastic extensions of EVM \citep{Acebes2015, Batselier2017} and Monte Carlo
simulation frameworks \citep{Kwak2007} have modelled performance uncertainty at
the project level, establishing that CPI and SPI follow empirically
characterisable distributions over project life. Empirical evidence on
contractor performance variability from EVM datasets
\citep{Acebes2015, Batselier2017} provides the calibration basis for the
stochastic efficiency parameter $\eta_{i,t}$ that parameterises the stochastic
program at Level~2 and the CMDP environment at Level~3.

S-curve analysis, formalised by \citet{Kenley2009} for construction cash flows,
established that budget consumption follows characteristic sigmoid profiles
whose shape parameters are project-type- and contract-form-dependent. The
regularized incomplete Beta function provides a flexible single-family
parameterisation of this profile \citep{barraza2004probabilistic}, and is
adopted directly in the planned expenditure model of the present formulation.

Large-scale empirical evidence on project performance is provided by
\citet{Merrow2011} through the Independent Project Analysis database, and by
\citet{Flyvbjerg2003, Flyvbjerg2018} through infrastructure cost-overrun
datasets spanning decades and geographies. These datasets inform both the
calibration of uncertainty distributions in the present model and the
benchmarks against which synthetic environment parameters are validated.
Complementary empirical benchmarks are drawn from AACE International and
Construction Industry Institute cost-engineering databases
\citep{christensen1993cost}.

Contract governance structures --- particularly FIDIC forms (Red, Yellow, and
Silver Books) --- prescribe the financial interaction between owner and
contractor through milestone payment schedules, advance payment and recovery
mechanisms, and retention clauses \citep{FIDIC2017}. The Estimate at
Completion $\mathrm{EAC}_{i,t} = \mathrm{ACWP}_{i,t} + (\mathrm{BAC}_i -
\mathrm{BCWP}_{i,t})/\mathrm{CPI}_{i,t}$ extends CPI into a forward-looking
cost trajectory signal; \citet{christensen1993cost} established across 155
contracts that this formulation produces the most accurate and stable
cost-at-completion forecasts, with reliability increasing beyond the 20\%
completion mark \citep{Lipke2003}. The present model adopts EAC as the
second termination condition signal, directly parallel to the forecast finish
deviation used in Condition~1. Under FIDIC
Sub-Clause~14.2, the advance payment is an interest-free mobilisation loan
recovered through percentage deductions from interim payment certificates: if
not otherwise specified in the Appendix to Tender, deductions commence once
cumulative certified interim payments exceed 10\% of the Accepted Contract
Amount and are applied at a rate of 25\% of each subsequent payment certificate
until the advance is fully recovered \citep{FIDIC2017}. The present model
parameterises this mechanism with a project-specific commencement threshold
$\psi_i$ and deduction rate $\delta_i^{\mathrm{rec}}$, accommodating the full
range of particular conditions encountered in practice. These mechanisms translate
directly into the cash inflow and outflow dynamics formulated in
Section~\ref{sec:problem_formulation}.

% [PLACEHOLDER: Additional empirical PM literature on contractor performance
%  variability (CPI/SPI distributions from Lipke 2003, Henderson 2004) to
%  be incorporated once final calibration data access is confirmed.]

% ------------------------------------------------------------
\subsection{Project Portfolio Management: Selection and Governance}
\label{subsec:lr_ppm}
% ------------------------------------------------------------

Project portfolio management formalises the problem of allocating an
organisation's finite resources across a competing set of candidate projects to
maximise strategic value \citep{Archer1999, Cooper2001}. The literature
distinguishes two nested decision layers: \emph{portfolio selection} --- the
binary or mixed-integer optimisation of which projects to include --- and
\emph{portfolio governance} --- the ongoing operational oversight of approved,
active projects.

Portfolio selection models have been extensively developed. Markowitz-style
mean--variance frameworks adapted from financial portfolio theory
\citep{Markowitz1952} have been applied to project baskets to trade off
expected NPV against variance \citep{Eilat2006}. Multi-criteria decision
analysis extensions incorporate strategic alignment, risk ratings, and resource
interdependencies \citep{Martinsuo2007, Killen2008}. Metaheuristic approaches
--- genetic algorithms, particle swarm optimisation --- address large
combinatorial instances \citep{Zhang2020}. Resource-constrained portfolio
selection under capital rationing has been studied by \citet{kavadias2003resource},
whose no-borrowing constraint is adopted directly in the present formulation.

The present work \emph{explicitly adopts the output of portfolio selection as
an exogenous input}. The selection problem is treated as solved upstream; the
focus is entirely on the operational governance layer that receives a committed
portfolio and must manage it through execution. This decomposition is consistent
with the two-stage governance model of \citet{Turner2014} and the PMI Standard
for Portfolio Management \citep{PMI2017}, which separates strategic selection
from operational execution oversight. In doing so, this work builds directly on
the selection literature while advancing the frontier into the governance domain
that has received substantially less formal analytical treatment.

% [PLACEHOLDER: Governance-layer literature on portfolio performance
%  monitoring and control (Dye and Pennypacker 1999, Levine 2005) to be
%  incorporated once screened corpus is finalised.]

% ------------------------------------------------------------
\subsection{PPM Budgeting, Cash-Flow Modelling, and Uncertainty}
\label{subsec:lr_cashflow}
% ------------------------------------------------------------

Within portfolio governance, budget allocation across active projects over time
represents a high-frequency, high-stakes operational decision. Despite its
practical centrality, this specific problem has received limited formal
treatment relative to portfolio selection. Existing work addresses project-level
cash-flow modelling \citep{Kenley2009, Park2004}, EVM-based forecasting of
individual project cost-at-completion \citep{Kim2002, Acebes2015}, and
multi-project resource levelling \citep{Lova2000, Peteghem2010}. However,
these contributions treat budgets as deterministic constraints rather than as
dynamic decision variables subject to portfolio-level sequential optimisation.

Cash-flow uncertainty in the present model arises from a single source:
the budget-to-progress efficiency $\eta_{i,t}$, which governs how much
physical progress a given budget allocation yields in each period
\citep{Acebes2015, Batselier2017}. Uncertain productivity determines whether
and when each milestone threshold is crossed, which in turn determines the
timing of the contractually fixed payment $R_{i,j}^{\mathrm{net}}$. Payment
delays between certification and receipt are not modelled as an independent
uncertainty source. This relaxation is defensible: the performance-driven
channel is the structurally primary source of cash-flow timing uncertainty
in milestone-payment contracts, and isolating it yields a cleaner
identification of the allocation problem without loss of the essential
uncertainty framing. Other sources of cash-flow uncertainty --- scope
changes, external market shocks \citep{Doloi2012} --- are treated as exogenous
and outside the scope of the present formulation. Stochastic programming
\citep{Birge2011} provides the canonical OR framework for optimisation under
such uncertainty, and its two-stage variant has been applied to project
scheduling under resource uncertainty \citep{Lova2000}. However, multi-stage
stochastic programs for portfolio-level sequential budget
allocation face scenario-tree explosion as the number of projects and
periods increases. The single uncertainty source $\eta_{i,t}$ now branches the
scenario tree at every period: if $\eta_{i,t}$ takes $K$ discrete values at
each of $H$ periods for each of $n$ projects, the tree produces $K^{nH}$ leaf
nodes. The effective planning horizon is itself portfolio-dependent---determined
by the latest project finish---compounding the explosion, since $H$ grows with
the most delayed project in the committed set. Both effects are formalised in
Section~\ref{subsec:level2}.

The gap in the literature is precise: \emph{no established framework models the
portfolio-level budget allocation decision as a sequential stochastic
optimisation problem in which a central decision-maker allocates funds across
multiple active projects whose cash inflow trajectories are uncertain through
stochastic budget-to-progress efficiency, subject to contractual milestone
certification windows that bound the speed at which progress converts to cash
inflows.} The present work
addresses this gap. The objective is pure discounted net cash flow; no deviation penalty enters
the objective function. This is a deliberate design choice: the consequences
of sustained deviation --- termination and loss of all future milestone
payments and retention --- are sufficient to make deviation undesirable
without a separate penalty term, and including one would introduce an
arbitrary scalarisation weight with no natural calibration from contract data
\citep{Ng1999}. Symmetrically, the discount factor $\gamma$ encodes not only
time preference but the opportunity cost of tied-up capital in a competitive
environment where future projects are foregone while current ones remain
incomplete \citep{Damodaran2012}. This interpretation gives $\gamma$ a
concrete calibration anchor beyond abstract preference --- it can be estimated
from the firm's observed pipeline redeployment rates or from prevailing
cost-of-capital benchmarks. The termination settlement further grounds the model in real asset values: the
settlement payment $R_i^{\mathrm{term}} = P_{i,t}^{\mathrm{actual}} \cdot
CP_i - (A_i + \sum_{j \in \mathcal{M}_i^{\mathrm{cert}}} R_{i,j}^{\mathrm{net}})$
computes the difference between the contractor's entitlement for work delivered
and cumulative inflows already received, which can be positive or negative
depending on progress at termination relative to certified payments
\citep{Dixit1994, Huchzermeier2001}, rather than an arbitrarily weighted
deviation penalty. A consequence is that the learned policy may
produce allocation trajectories that concentrate expenditure around milestone
certification thresholds, mirroring documented contractor financial strategy
in milestone-payment contracts \citep{Kenley2009, Kaka1993, Arditi2005}.

% [PLACEHOLDER: Cash-flow modelling literature specific to EPC sector
%  (Kaka and Price 1993, Navon 1996) to be incorporated. Also: literature
%  on advance payment and retention mechanics in international contracting
%  (Eggers 2019, CIDB 2015) to be cross-checked against screened corpus.]

% ------------------------------------------------------------
\subsection{Contractor Underperformance, Cure Periods, and Termination}
\label{subsec:lr_interventions}
% ------------------------------------------------------------

When a contractor fails to meet schedule or cost performance thresholds,
international EPC contracts prescribe a structured escalation sequence before
termination may be exercised. Under FIDIC Sub-Clause~15.1, the client must
issue a formal notice of default and allow a cure period --- typically 28~days
--- during which the contractor is obliged to demonstrate measurable recovery
before the client may proceed to termination \citep{FIDIC2017, Ndekugri2008}.
This cure period is not a discretionary grace; it is a contractual obligation
whose omission constitutes wrongful termination. Empirical studies of
international EPC disputes confirm that failure to observe the cure period is
among the most frequently litigated grounds in contractor termination claims
\citep{Ndekugri2008, Ennis2021}.

The behavioural project management literature has studied the conditions under
which organisations fail to terminate underperforming projects at the
economically rational point. Escalation of commitment \citep{Staw1987} ---
the tendency to continue investing in failing projects beyond what expected
value justifies --- has been documented across construction, IT, and
infrastructure programmes. Real options theory \citep{Dixit1994,
Huchzermeier2001} provides the normative counterpart: termination is the
exercise of an abandonment option whose value exceeds continuation NPV once
performance has deteriorated past a recoverable threshold. Together these
streams establish that the decision to persist with or terminate an
underperforming project is neither purely contractual nor purely financial ---
it involves performance signals, contractual obligations, and timing all at
once.

The present model integrates all three dimensions into a single environment
mechanic governed by two conditions. A \emph{progress delay breach} fires
whenever the lag between planned and actual progress exceeds the contractual
tolerance $\delta^{\max}_i$, i.e.\ $\delta_{i,t} = PV_i(\tau_i) - EV_{i,t} >
\delta^{\max}_i$, grounded in the FIDIC engineer's authority to issue a notice
when progress falls materially behind programme \citep{FIDIC2017}. A
\emph{hard deadline breach} fires unconditionally at episode period $t \geq
f_i + \Delta^{\max}_i$, where $\Delta^{\max}_i$ is the contractually permitted
finish overrun, and triggers immediate termination regardless of the
cure-period counter state. When a progress delay breach is active, the
cure-period counter $\theta^{\mathrm{rem}}_{i,t}$ decrements by one; it
resets to $\theta_i$ only when no breach is present. Termination executes when
the counter reaches zero. This mechanic is grounded
simultaneously in FIDIC contract law \citep{FIDIC2017, Ndekugri2008}, the real
options framing of the abandonment threshold \citep{Dixit1994,
Huchzermeier2001}, the EVM literature's validation of EAC as a stable
forward cost signal \citep{christensen1993est, Lipke2003}, and the
escalation-of-commitment literature's evidence that observable performance
signals must be sustained --- not transient --- before termination is
justified \citep{Staw1987}. The counter $\theta^{\mathrm{rem}}_{i,t}$ is part of the observable portfolio
state, giving the RL agent an anticipatory signal that reflects contractual
reality without granting it discretionary control over the termination outcome.
From the portfolio manager's perspective, this signal serves a dual function:
it flags projects approaching contractual failure while also marking the point
at which continued allocation yields diminishing returns relative to the
settlement value of work already delivered. A rational manager reading a
counter approaching zero in a high-opportunity-cost environment would recognise
that the real value of future milestone payments is eroding while the
settlement value of work in place is fixed in nominal terms --- creating a
natural reallocation incentive that the model captures through the interaction
of $\gamma$ and $S_i = EV_{i,t} \cdot P_i - \mathrm{CF}^{+}_{i}$ without
requiring explicit inflation modelling \citep{Dixit1994, FIDIC2017}.


% ------------------------------------------------------------
\subsection{Reinforcement Learning for Sequential Resource Allocation}
\label{subsec:lr_rl}
% ------------------------------------------------------------

Reinforcement learning provides a principled framework for sequential
decision-making under uncertainty \citep{Sutton2018}. The MDP
formalisation --- state, action, transition, reward --- maps naturally onto the
portfolio budget allocation problem, and the constrained MDP extension
\citep{altman1999constrained} accommodates the hard budget feasibility constraints that
the portfolio manager must satisfy at every period.

RL has been applied to project scheduling \citep{Peng2021, Liu2023},
resource-constrained project scheduling problems (RCPSP) \citep{Tassel2023},
and supply chain inventory management \citep{Giannoccaro2002}. In each domain,
 RL demonstrates advantages over classical OR methods when
state spaces are large, transition dynamics are stochastic, the problem is
non-stationary over the decision horizon, and the horizon length is itself
uncertain at the outset of planning. These are precisely the characteristics of the
portfolio budget allocation problem at Levels~2 and~3.

Proximal Policy Optimisation \citep{Schulman2017} is adopted as the solution
algorithm for its stability in continuous action spaces, its sample efficiency
relative to earlier policy-gradient methods, and its established use in
operations management domains. Domain randomisation \citep{Tobin2017} ---
randomising environment parameters during training to improve out-of-sample
policy robustness --- is adopted to ensure the learned policy generalises across
the range of project types, contract structures, and market conditions
represented in the calibration data. Walk-forward validation
\citep{Cerqueira2020} is used as the empirical evaluation protocol to guard
against data leakage given the time-series structure of project cash-flow data.

Direct application of RL to portfolio-level budget allocation remains absent
from the literature. The present work provides the first formulation of this
problem as a CMDP and the first empirical demonstration that a PPO agent can
learn non-trivial allocation policies in a calibrated synthetic environment.

% ------------------------------------------------------------
\subsection{Synthesis: Positioning the Present Contribution}
\label{subsec:lr_synthesis}
% ------------------------------------------------------------

Table~\ref{tab:litmap} maps representative works across the four literature
streams against the key modelling dimensions of the present framework,
as generated by the automated synthesis pipeline described in
Section~\ref{subsubsec:prisma}. Three structural gaps emerge from this mapping.

\textbf{Gap~1 --- Selection versus operational governance.}\;
The transition from portfolio \emph{selection} to portfolio \emph{operational
governance} has not been modelled as a sequential decision problem. Selection
literature assumes away post-selection management; project scheduling literature
ignores portfolio interdependencies and the shared capital constraint. No
existing work closes both simultaneously.

\textbf{Gap~2 --- Project-level versus portfolio-level uncertainty propagation.}\;
Existing stochastic models of project cash-flow uncertainty operate at the
individual project level and do not propagate uncertainty upward into a
portfolio-level decision architecture governed by a shared budget constraint.
The feedback loop between a project's stochastic efficiency trajectory and
its portfolio-level funding entitlement remains unmodelled.

\textbf{Gap~3 --- Absence of endogenous termination mechanics.}\;
The joint modelling of a progress-delay breach condition tracked against a
contractual tolerance, a decrementing cure-period counter governing
grace-period exhaustion, an independent hard-deadline breach condition that
triggers immediate termination regardless of the cure-period state, and a
termination settlement cash flow reconciling earned value against
cumulative payments received, has not been embedded within a
portfolio-level sequential decision framework in which budget allocation
is the sole control variable.

\medskip
These gaps are not isolated oversights; they reflect an incomplete progression
along the path toward a deployable portfolio-level decision-support system.
Figure~\ref{fig:maturity} situates this progression as three sequential steps.

\textbf{Step~1} --- domain grounding --- is well established by the literature
reviewed above: contract mechanics, EVM performance signals, S-curve expenditure
profiles, and portfolio selection models collectively provide the inputs and
boundary conditions that any operational governance model requires.

\textbf{Step~2} --- tractable problem formulation --- is the gap this paper
closes. The contribution is a verified, configurable CMDP environment that
encodes the full FIDIC cash-flow cycle and a two-condition termination mechanic,
released as a Gymnasium-compatible package with a complete SQLite audit ledger.
Baseline greedy heuristics are benchmarked against a full-information optimum,
producing a quantified optimality gap that grows monotonically with portfolio
complexity and establishes the performance target for future learned agents.

\textbf{Step~3} --- validated deployment within incumbent project management
information systems such as Oracle Primavera~P6 --- remains future work,
contingent on the trained agents that Step~2 now makes possible.

\begin{figure}[htbp]
\centering
\begin{tikzpicture}[
    node distance  = 0cm,
    stepbox/.style = {rectangle, draw, rounded corners=3pt,
                      minimum width=3.6cm, minimum height=2.8cm,
                      align=center, font=\small, text width=3.2cm},
    arrow/.style   = {->, >=stealth, thick},
    label/.style   = {font=\scriptsize, align=center, text width=3.2cm}
]

\node[stepbox, fill=gray!15] (s1) at (0,0)
    {\textbf{Step 1}\\[4pt] Domain Grounding\\[6pt]
     \scriptsize Contract mechanics\\EVM signals\\S-curve profiles\\
     Portfolio selection};

\node[stepbox, fill=gray!8, right=0.6cm of s1] (s2)
    {\textbf{Step 2} \\[1pt] \emph{(this work)}\\[4pt]
     Problem Formulation\\[6pt]
     \scriptsize CMDP environment\\FIDIC cash-flow\\
     Termination mechanic\\Gymnasium API};

\node[stepbox, draw=gray, dashed, right=0.6cm of s2] (s3)
    {\textbf{Step 3}\\[4pt] Validated Deployment\\[6pt]
     \scriptsize RL agent training\\OR benchmark\\
     Primavera~P6 integration\\[4pt]\emph{(future work)}};

\draw[arrow] (s1) -- (s2);
\draw[arrow, dashed, color=gray] (s2) -- (s3);

\end{tikzpicture}
\caption{Maturity progression toward a deployable portfolio-level
decision-support system. Step~2, completed by the present work, enables
Step~3. Dashed border and arrow denote future work.}
\label{fig:maturity}
\end{figure}

\medskip
The present work addresses all three gaps and completes Step~2 of this
progression. The Gymnasium-compatible CMDP environment --- with
parameterisable contract mechanics, stochastic budget-to-progress efficiency
as the single uncertainty source, and an endogenous termination mechanic ---
constitutes the primary methodological contribution. Baseline heuristics are
evaluated against a full-information optimum in Section~\ref{sec:evaluation},
providing a quantified optimality gap that grows monotonically with portfolio
complexity and defines the improvement target for future RL agents.


\begin{table}[htbp]
    \centering
    \caption{Representative literature mapped against modelling dimensions of the present framework. \textit{Decision type}: S\,=\,selection, Sc\,=\,scheduling, B\,=\,budgeting. \checkmark~=~fully addressed;\; $\circ$~=~partially addressed;\; {--}~=~not addressed.}
    \label{tab:litmap}
    \setlength{\tabcolsep}{6pt}
    \renewcommand{\arraystretch}{1.35}
    \rowcolors{3}{rowodd}{white}
    \resizebox{\textwidth}{!}{%
    \begin{tabular}{>{\raggedright\arraybackslash}p{3.8cm}cc>{\columncolor{rarebg}}cc>{\columncolor{rarebg}}c>{\columncolor{rarebg}}c}
    \toprule
    \rowcolor{headbg}\color{headfg}\textbf{Work} & \color{headfg}\textbf{Portfolio} & \color{headfg}\textbf{Decision} & \cellcolor{headbg}\color{headfg}\textbf{Sequential} & \color{headfg}\textbf{Stochastic} & \cellcolor{headbg}\color{headfg}\textbf{Contract} & \cellcolor{headbg}\color{headfg}\textbf{Termination} \\
    \rowcolor{headbg}\color{headfg}\textbf{} & \color{headfg}\textbf{level} & \color{headfg}\textbf{type} & \cellcolor{headbg}\color{headfg}\textbf{decisions} & \color{headfg}\textbf{cash flows} & \cellcolor{headbg}\color{headfg}\textbf{mechanics} & \cellcolor{headbg}\color{headfg}\textbf{settlement} \\
    \midrule
    \citet{Yamakawa2019} & \textcolor{symgreen}{\checkmark} & \textsc{\small S} & \textcolor{symgrey}{--} & \textcolor{symamber}{$\circ$} & \textcolor{symgrey}{--} & \textcolor{symgrey}{--} \\
    \citet{KOSZTYAN2022} & \textcolor{symgreen}{\checkmark} & \textsc{\small S} & \textcolor{symgrey}{--} & \textcolor{symamber}{$\circ$} & \textcolor{symgrey}{--} & \textcolor{symgrey}{--} \\
    \citet{touran2006mathematical} & \textcolor{symgreen}{\checkmark} & \textsc{\small Sc/B} & \textcolor{symgrey}{--} & \textcolor{symamber}{$\circ$} & \textcolor{symgrey}{--} & \textcolor{symgrey}{--} \\
    \citet{touran2010probabilistic} & \textcolor{symgreen}{\checkmark} & \textsc{\small Sc/B} & \textcolor{symgrey}{--} & \textcolor{symamber}{$\circ$} & \textcolor{symgrey}{--} & \textcolor{symgrey}{--} \\
    \citet{GHAPANCHI2012791} & \textcolor{symgreen}{\checkmark} & \textsc{\small S} & \textcolor{symgrey}{--} & \textcolor{symamber}{$\circ$} & \textcolor{symgrey}{--} & \textcolor{symgrey}{--} \\
    \citet{pourahmadi2015scenario} & \textcolor{symgreen}{\checkmark} & \textsc{\small S} & \textcolor{symgrey}{--} & \textcolor{symamber}{$\circ$} & \textcolor{symgrey}{--} & \textcolor{symgrey}{--} \\
    \citet{elsabek2018coordination} & \textcolor{symgrey}{--} & \textsc{\small B} & \textcolor{symgrey}{--} & \textcolor{symgrey}{--} & \textcolor{symgreen}{\checkmark} & \textcolor{symgrey}{--} \\
    \citet{beseiso2021} & \textcolor{symgreen}{\checkmark} & \textsc{\small S} & \textcolor{symgrey}{--} & \textcolor{symamber}{$\circ$} & \textcolor{symgrey}{--} & \textcolor{symgrey}{--} \\
    \citet{karimi2021scenario} & \textcolor{symgreen}{\checkmark} & \textsc{\small S} & \textcolor{symgrey}{--} & \textcolor{symamber}{$\circ$} & \textcolor{symgrey}{--} & \textcolor{symgrey}{--} \\
    \citet{bahroun2023artificial} & \textcolor{symgrey}{--} & \textsc{\small Sc} & \textcolor{symgreen}{\checkmark} & \textcolor{symgrey}{--} & \textcolor{symgrey}{--} & \textcolor{symgrey}{--} \\
    \citet{RANJBAR2022116134} & \textcolor{symgreen}{\checkmark} & \textsc{\small S/Sc} & \textcolor{symgrey}{--} & \textcolor{symgrey}{--} & \textcolor{symamber}{$\circ$} & \textcolor{symgrey}{--} \\
    \citet{szwarcfiter2023balancing} & \textcolor{symgreen}{\checkmark} & \textsc{\small Sc} & \textcolor{symgreen}{\checkmark} & \textcolor{symamber}{$\circ$} & \textcolor{symgrey}{--} & \textcolor{symgrey}{--} \\
    \citet{asghari2023} & \textcolor{symgrey}{--} & \textcolor{symgrey}{--} & \textcolor{symgreen}{\checkmark} & \textcolor{symamber}{$\circ$} & \textcolor{symgrey}{--} & \textcolor{symgrey}{--} \\
    \citet{mazhari2024} & \textcolor{symgreen}{\checkmark} & \textsc{\small S/B} & \textcolor{symgreen}{\checkmark} & \textcolor{symgrey}{--} & \textcolor{symgrey}{--} & \textcolor{symgrey}{--} \\
    \citet{JIANG2024101016} & \textcolor{symgreen}{\checkmark} & \textcolor{symgrey}{--} & \textcolor{symgreen}{\checkmark} & \textcolor{symamber}{$\circ$} & \textcolor{symgrey}{--} & \textcolor{symgrey}{--} \\
    \citet{Nabot2025} & \textcolor{symgreen}{\checkmark} & \textsc{\small Sc} & \textcolor{symgreen}{\checkmark} & \textcolor{symgrey}{--} & \textcolor{symgreen}{\checkmark} & \textcolor{symgreen}{\checkmark} \\
    \citet{Darvish2025} & \textcolor{symgreen}{\checkmark} & \textsc{\small B} & \textcolor{symgreen}{\checkmark} & \textcolor{symamber}{$\circ$} & \textcolor{symamber}{$\circ$} & \textcolor{symgreen}{\checkmark} \\
    \citet{Gholami2025} & \textcolor{symgreen}{\checkmark} & \textsc{\small Sc} & \textcolor{symgreen}{\checkmark} & \textcolor{symamber}{$\circ$} & \textcolor{symgrey}{--} & \textcolor{symgreen}{\checkmark} \\
    \citet{Kehinde2025} & \textcolor{symgreen}{\checkmark} & \textsc{\small B} & \textcolor{symgreen}{\checkmark} & \textcolor{symgreen}{\checkmark} & \textcolor{symgrey}{--} & \textcolor{symgrey}{--} \\
    \citet{Ajibade2025} & \textcolor{symgreen}{\checkmark} & \textsc{\small Sc} & \textcolor{symgreen}{\checkmark} & \textcolor{symamber}{$\circ$} & \textcolor{symgrey}{--} & \textcolor{symgrey}{--} \\
    \citet{Huddar2025} & \textcolor{symgreen}{\checkmark} & \textsc{\small B} & \textcolor{symgreen}{\checkmark} & \textcolor{symamber}{$\circ$} & \textcolor{symgrey}{--} & \textcolor{symgrey}{--} \\
    \citet{Akbari2026} & \textcolor{symgreen}{\checkmark} & \textsc{\small B} & \textcolor{symgreen}{\checkmark} & \textcolor{symamber}{$\circ$} & \textcolor{symgreen}{\checkmark} & \textcolor{symgreen}{\checkmark} \\
    \citet{ebnerasoul2023} & \textcolor{symgreen}{\checkmark} & \textsc{\small S} & \textcolor{symgrey}{--} & \textcolor{symamber}{$\circ$} & \textcolor{symgrey}{--} & \textcolor{symgrey}{--} \\
    \citet{noroozi2016probabilistic} & \textcolor{symgreen}{\checkmark} & \textsc{\small B} & \textcolor{symgrey}{--} & \textcolor{symamber}{$\circ$} & \textcolor{symgrey}{--} & \textcolor{symgrey}{--} \\
    \citet{JIANG2023104817} & \textcolor{symgrey}{--} & \textsc{\small B} & \textcolor{symgreen}{\checkmark} & \textcolor{symgrey}{--} & \textcolor{symgrey}{--} & \textcolor{symgrey}{--} \\
    \cmidrule{1-7}
    \rowcolor{pwbg}\color{headfg}\textbf{Present work} & \cellcolor{pwbg}\color{headfg}\checkmark & \color{headfg}\textbf{B} & \cellcolor{pwbg}\color{headfg}\checkmark & \color{headfg}\checkmark & \cellcolor{pwbg}\color{headfg}\checkmark & \cellcolor{pwbg}\color{headfg}\checkmark \\
    \bottomrule
    \end{tabular}%
    }
  \end{table}